\par\Needspace{6\baselineskip}\medskip
\hypertarget{native-111-article082}{}\pdfbookmark[2]{82.}{article-082}
\begin{center}82.\end{center}
Nehmen wir schließlich an, es seien drei Radiusvektoren $r$, $r'$, $r''$ gegeben, die mit einer beliebigen, von der Sonne aus in der Bahnebene gezogenen Geraden die Winkel $N$, $N'$, $N''$ bilden. Behält man die übrigen Bezeichnungen bei, so erhält man die Gleichungen
\[
\begin{aligned}
&\frac{p}{r}=1+e\cos(N-\Pi)\\[4pt]
\mathrm{I.}\qquad&\frac{p}{r'}=1+e\cos(N'-\Pi)\\[4pt]
&\frac{p}{r''}=1+e\cos(N''-\Pi),
\end{aligned}
\]
aus denen sich $p$, $\Pi$ und $e$ auf verschiedene Arten ermitteln lassen. Will man $p$ vor den übrigen Größen berechnen, so multipliziert man die drei Gleichungen I der Reihe nach mit $\sin(N''-N')$, $-\sin(N''-N)$ und $\sin(N'-N)$. Durch Addition der Produkte ergibt sich nach Lemma I aus Art.\ 78
\[
p=\frac{\sin(N''-N')-\sin(N''-N)+\sin(N'-N)}{\dfrac1r\sin(N''-N')-\dfrac1{r'}\sin(N''-N)+\dfrac1{r''}\sin(N'-N)}\mathpunct{\cdot}
\]
Dieser Ausdruck verdient eine nähere Betrachtung. Der Zähler geht offenbar über in
\[
\begin{aligned}
&=2\sin\tfrac12(N''-N')\cos\tfrac12(N''-N')\\
&\qquad-2\sin\tfrac12(N''-N')\cos(\tfrac12N''+\tfrac12N'-N)\\
&=4\sin\tfrac12(N''-N')\sin\tfrac12(N''-N)\sin\tfrac12(N'-N).
\end{aligned}
\]
Setzt man ferner
\[
\begin{aligned}
r'r''\sin(N''-N')&=n\\
rr''\sin(N''-N)&=n'\\
rr'\sin(N'-N)&=n'',
\end{aligned}
\]
so ist klar, dass $\tfrac12n$, $\tfrac12n'$, $\tfrac12n''$ die Flächeninhalte der Dreiecke zwischen dem zweiten und dritten, dem ersten und dritten beziehungsweise dem ersten und zweiten Radiusvektor sind. Daraus erkennt man leicht, dass in der neuen Formel
\SourceLeaf{112}
\[
p=\frac{4\sin\tfrac12(N''-N')\sin\tfrac12(N''-N)\sin\tfrac12(N'-N)\mathbin{.}rr'r''}{n-n'+n''}
\]
der Nenner das Doppelte des Flächeninhalts des Dreiecks ist, dessen Ecken die Endpunkte der drei Radiusvektoren sind, also die drei Orte des Himmelskörpers im Raum. Liegen diese Orte nahe beieinander, so ist dieser Flächeninhalt stets sehr klein, und zwar von dritter Ordnung, wenn man $N'-N$, $N''-N'$ als kleine Größen erster Ordnung auffasst. Zugleich folgt daraus, dass Fehler in einer oder mehreren der Größen $r$, $r'$, $r''$, $N$, $N'$, $N''$, so gering sie auch sein mögen, einen sehr großen Fehler bei der Bestimmung von $p$ verursachen können. Dieses Verfahren zur Ermittlung der Abmessungen der Bahn kann daher nur dann eine hohe Genauigkeit liefern, wenn die drei heliozentrischen Orte beträchtliche Abstände voneinander haben.

Sobald der Halbparameter $p$ gefunden ist, bestimmt man $e$ und $\Pi$ nach dem Verfahren aus Art.\ 79, indem man zwei beliebige der Gleichungen I miteinander kombiniert.
